% tex/sections/outlook.tex
% ──────────────────────────────────────────────────────────────────────────────
\chapter{Outlook}
\label{chap:outlook}
% ──────────────────────────────────────────────────────────────────────────────
%
% TODO: Write outlook.
%
% Suggested structure:
%   - Future work directly motivated by this thesis
%   - Broader research directions opened up
%   - Proposed experiments or improvements to methodology

The results of this thesis establish a baseline for DAS signatures of sub-Rayleigh and Supershear crack propagation in a simplified geometry, the following additions would improve the physical realism and practical applicability of the model. Based on the results and methods presented in this thesis, there are many open lines of inquiry. Because the model described in the simple case is a basic two layer model, a next step could be to add multiple layers, especially because cracks in snow usually propagate in a weak layer \parencite{schweizerSnowAvalancheFormation2003}. While this was done partially in the MPM-guided simulations, where a 5-layer model was implemented, the snow slab was still approximated as a homogeneous material with the same model properties above and below the weak layer, which is a major simplification. Here, there is also possibility to run the simulations described in this thesis for different snow parameters (density, $v_p$ and $v_s$, elastic modulus), because \cite{barteltPhysicalSNOWPACKModel2002,mulakNumericalInvestigationMixedmode2019,bobillierMicromechanicalModelingSnow2020} have all shown that snow properties not only change over time, but also have an influence on crack propagation. This would be useful to better understand DAS data from different times of the year or day, when there is the possibility of snow properties having changed. \par

In both simulation types run, there is also the possibility of adding topography or a slope angle to the simulation, because as indicated by \cite{anceySnowAvalanches2001,trottetTransitionSubRayleighAnticrack2022}, avalanches are gravity driven and the crack propagation regimes can change greatly with the presence of topography, which also has an effect on the seismic signatures. This would also help identify how the signatures observed here change when the DAS cable is not on flat ground. All of the results shown in this thesis assume perfect coupling, which is not always the case in the field, especially when topography is present, this is another reason why adding topography is important. \par 

Because of the possibility for transition between the different propagation regimes, meaning that the propagation speed changes, as described by \cite{trottetTransitionSubRayleighAnticrack2022}, which is not only driven by topography and slope angle, the constant propagation speed prescribed in the simulations in this thesis are a major limitation. Another next step would be to have source propagation accelerate with increasing propagation distance, instead of having a constant propagation speed, which is not very plausible on a mountain slope. The MPM model which was used to obtain the moment tensors in this case had a slope angle of 45 degrees \parencite{trottetTransitionSubRayleighAnticrack2022}. This means that there is no transition between modes in the data itself, but using these moment tensors obtained from data at a slope in flat terrain in SALVUS is also a major simplification. \par 

Furthermore, the crack modeled in this thesis is linear and at constant depth. The moment tensor extraction from MPM used here averages the stress field over the weak layer thickness, which loses information about the vertical structure of the source. MPM simulations of weak layer cracking and collapse have shown that this is not the case \parencite{trottetTransitionSubRayleighAnticrack2022}, so a next step would be to make the crack propagation modeled by seismic sources more physical. This could be done by having "jumps" of the crack into different depths of the snowpack or having multiple lines of sources firing at the same time to represent multiple smaller cracks. In the MPM case, a more refined approach would be extracting the moment tensor at individual MPM particle positions without averaging, preserving the depth-dependent stress distribution and potentially revealing differences in the Lamb wave mode excitation that are hidden by the averaging. This would also allow the crack depth to vary with propagation distance. \par

Another limitation is that the simple and CZM simulations cannot account for the physical collapse
of the weak layer that characterizes anticrack propagation
\parencite{heierliAnticrackNucleationTriggering2008,
heierliAnticracksNewTheory2003}.This is because spectral element methods solve only the elastic wave equation in the surrounding medium and cannot
represent inelastic failure of material elements without explicit element
removal \parencite{podolskiyReviewFiniteelementModelling2013}. The
prescribed moment tensor source in these models approximates the
mechanical effect of the collapse but does not model the collapse
kinematics themselves. In the MPM-guided simulations, this limitation is
partially addressed: the inelastic collapse process is modelled fully
within MPM, and its mechanical effect on the surrounding elastic medium
is transferred to SALVUS via the extracted moment tensor time series.
However, the elastic wave field generated by the collapse in the MPM
simulation is not re-injected into SALVUS, only the quasi-static
stress redistribution is captured through the moment tensors. Dynamic
wave radiation generated within the MPM domain itself is therefore not
included in the SALVUS simulation. The presence of these signals could also change the seismic signal. This could be verified by field tests of weak layer collapse measured by a DAS cable. \par 

The quasi-static strain field identified in the MPM results could be further investigated as an early warning signal. The amplitude of this field relative to the DAS noise floor determines whether it would be detectable in a field deployment, and this depends on the specific DAS instrument, gauge length, and environmental noise. A field comparison using a PST experiment instrumented with both DAS and geophones would allow the quasi-static component to be identified in real data, validating the MPM simulation prediction. Such field experiments have been conducted for dynamic wave signatures
\parencite{bergfeldSignaturesSubRayleighSupershear2025,herwijnenSeismicSensorArray2011} but the quasi-static DAS signature has not yet been specifically targeted. \par

In addition to the limitations described, the MPM data used to obtain the moment tensors resulted in a sub-Rayleigh rupture. While the results confirm the signatures observed in the simple model case, such validation was not obtained for the Supershear results or for a transitional regime. This would require MPM simulation data of those regimes, which was not available in this thesis. Here, there is a lot of potential for further validation of the signatures observed, because although they seem like useful precursors and identifiers in this data, more validation is needed. \par 

The last step in making the model described here physically accurate is to implement the 2D model described here in 3D, which would help to understand slope DAS data even better. Such a 3D model could represent data recorded further away from the crack, in this thesis, the 2D geometry constrains the cable to be parallel and or co-planar to the crack, whereas in the field, cables may be oriented at arbitrary angles to the crack propagation direction. It would also be very interesting to compare different cable orientations relative to the crack to see if the wavefield data recorded changes. \par 

